% ============================================================
% Part III — Desktop Section — M7 (Implications and Close) — DRAFT v1
% Mode: Stage 3 guided drafting. Plan APPROVED 2026-07-28:
%   M7.1 = 2 ¶¶ for the empirical program (gesture scale per
%   future-work rule — gestures, not plans); M7.2 = 1 ¶ for the
%   dissertation, SECTION SCALE ONLY per user ruling (O-11 deferred;
%   if O-11 resolves toward in-section conclusion, M7.2 is the
%   EXPANSION SITE — do not treat as dissertation-level close).
% RULINGS BANKED 2026-07-28:
%   (1) M7.2 register = section-scale close + O-11 expansion marker
%       (recommendation approved).
%   (2) Task-lever expansion = KEY QUOTATIONS woven into M5.1's
%       task-lever treatment (NOT a lit review, NOT an M1 supplement
%       pass) — queued as post-M7 pass on M5.1.
% STALE OUTLINE POINTERS corrected in rendering: "cost-ranked remedies
%   ladder" → ordered by the shared time each asks (cost-metaphor ban +
%   ladder removal); M7.2's "stall, flight, and unconvened WE" →
%   silence at the act / furnished flight / unconvened WE (no-stall
%   doctrine — the chain runs always).
% HARMONIZATION FLAG (assembly): the three-state proposal has a
%   gesture-scale rendering in III-5-Design-DRAFT-v1.tex ("restless
%   search, occupied-but-hollow, withdrawn"; surface/depth disagreement
%   as diagnostic; sleep-exit terminal). M7.1 ¶1 hands forward WITHOUT
%   reusing that wording; reconcile the two gestures once arc decision
%   (#1) resolves which section carries the full statement.
% LIT-SLOT (verified 2026-07-16, both sides): FCM-phenomenology ×
%   instrumentation = UNBUILT combination — 0/15 hard-channel units
%   carry FCM bridge edges; none of 34 VR-pedagogy units combines the
%   phenomenological boredom apparatus with instrumentation; nearest =
%   EEG-workload "boredom"-band labels (Makransky, Terkildsen & Mayer
%   2019; Parong & Mayer 2021) + single self-report items; DA-10
%   corroborates at network level (no citation traffic joins halves).
% ALL standing sweeps in force (And / record / pole / ladder / stall /
%   cost-metaphors / bare-incorporation / vague-frames).
% VERIFICATION LEDGER (index-entry-first per standing rule):
%   ped-sec-10 (Makransky, Terkildsen & Mayer 2019): EEG workload from
%     9-ch Advanced Brain Monitoring system, bands boredom/optimal/
%     overload, entry loci pp. 226–27 + 229–30 → cite 229–30.
%   ped-sec-13 (Parong & Mayer 2021; online-first 2020 — cite VoR
%     2021): boredom = instrument's own label for lowest workload band
%     (FAITHFUL, 234–36); IVR sig. higher boredom-band proportion +
%     lower EEG high-engagement (FAITHFUL, 234–36); presence d=1.74,
%     enjoyment d=1.20 both IVR>PPT sig. (FAITHFUL, 233–36); transfer
%     PPT>IVR sig. d=−0.71, retention marginal p=.058 (FAITHFUL, 233)
%     → "learned less by the study's transfer measure" (retention held
%     to marginal — NOT claimed). Band thresholds NOT printed (entry
%     shows ≤.40 vs <.40 discrepancy — avoid until PDF-pinned).
%   MLA: Parong & Mayer short title required IF ped-sec-08 (2018) also
%     enters Works Cited — normalize at assembly. Works Cited adds
%     queued: Makransky/Terkildsen/Mayer 2019 + Parong/Mayer 2021.
% M7.1 COMPLETE 2026-07-28 (2 ¶¶): ¶2 approved first pass; ¶1 v2
%   approved (concrete opener + Parong/Makransky specifics — the
%   field's own instruments already produced the surface/depth
%   divergence, read only as misallocated workload).
% M7.2 (1 ¶) — SECTION-SCALE CLOSE per user ruling 2026-07-28.
%   O-11 EXPANSION MARKER: if O-11 resolves toward an in-section
%   dissertation-level conclusion, THIS PARAGRAPH is the expansion
%   site; as drafted it closes the section's own argument and hands
%   the frame back as a finding — do not read it as the dissertation's
%   conclusion. DRAFTED 2026-07-28; v2 rulings applied ("demonstrated"
%   not "prove" — too absolute; silence clause anchored [assigned
%   analysis completable entirely outside; inside, only watching
%   asked]; GREATEST-DESIRE DOCTRINE: never "better route" — the chain
%   shows the person following the greatest desire, stronger outweighs
%   weaker [banked in feedback-chain-node-analysis-fundamental]; rest
%   approved) — revised spans APPROVED 2026-07-28. M7.2 COMPLETE.
% M7 MODULE COMPLETE 2026-07-28 (M7.1 2 ¶¶ · M7.2 1 ¶) — DESKTOP
%   SECTION DRAFTING COMPLETE M0–M7.
% ============================================================

% --- M7.1 For the empirical program (2 ¶¶, gesture scale) ---

The virtual-learning literature already measures something it calls
boredom, and the shape of its instrument shows exactly what is
missing. In the field's most instrumented media-comparison
experiments, boredom is the name of the lowest band of an EEG workload
metric: whatever falls beneath the scale's lower threshold is
classified as ``boredom,'' so that the word means nothing more than
workload's absence---the bottom of a single axis (Makransky,
Terkildsen and Mayer 229--30; Parong and Mayer, ``Cognitive and
Affective Processes'' 234--36). One of those experiments shows why the
label is not enough. In Parong and Mayer's media comparison, the
students in the headset reported the highest presence and the highest
enjoyment in the study, with large effects on both; the same students
were classified into the boredom band at a significantly higher rate
than the slideshow group, and they learned less by the study's
transfer measure (``Cognitive and Affective Processes'' 233--36). The
self-report channel found those students full; the physiological
channel found them empty; the outcome measure sided with the empty
reading. That is the divergence this section could only hear as an
echo in a student's prose, produced there by independent instruments
in a single published experiment---and the study can read it only as
misallocated workload, because boredom-as-low-workload has no
structure: it cannot distinguish the student searching for something
to attend to from the student occupied but held by nothing from the
student withdrawn from the world altogether. Heidegger's account
distinguishes exactly these, and no published study joins it to an
instrumented reading of learners inside a virtual world. The
combination is unbuilt, and the laboratory study this Part turns to
next was designed inside the frame that makes it buildable: a surface
reading of what occupies the student run concurrently with a depth
reading of what holds them, the disagreement between the channels
serving as the diagnostic.

The unconvened WE leaves the empirical program a narrower question,
and it comes with its own apparatus attached: what is the least shared
time that convenes the WE? The design section's four remedies already
stand in the order the question requires, by how much shared time each
asks of a group---the demonstration asks none, the traces ask none,
the scheduling assignment asks one hour of the group's own choosing,
the minutes make that hour accountable---so the redesigned deployment
is itself the instrument: introduce the remedies in that order, term
by term, and watch for the increment at which students begin to meet.
The collected data has already fixed the baseline any such observation
starts from---thirteen students asking for the meeting, one confirmed
convening---and the question's answer would be a finding about
\textit{chronos} as a design material: how small a reserve of shared
time suffices. Where the two proposals meet, a further question
becomes askable for the first time. This section marked as speculation
the thought that the convened WE might hold a student where the
solitary world could not; a deployment that both convenes the meeting
and listens on two channels is exactly the apparatus in which that
speculation would stop being one.

% --- M7.2 For the dissertation (1 ¶, section-scale close) ---

The desktop clinic completes a test the dissertation's frame has been
running since Part I: whether the vocabulary built there and
demonstrated in Part II against the richest of authored worlds---the chain from
perception to act, rhetorical incorporation, world-disclosedness and
co-disclosedness, the demand of veri(dis)similitude---could hold in
the sparest: procedure footage, plain rooms, proximity voice, deployed
under mandate to two hundred forty-one students. Everything those
students met was authored---the rooms, the footage, the hail spoken
before the world first rendered---and none of that authorship
determined an outcome. The chain ran always; each student's acts
completed where that student's own judgment routed them, most often in
the alternative the course itself had furnished. What the frame
contributes is the ability to say precisely what stood in the gap
between total authorship and undetermined outcome. Without it, the
survey's answers read as a feature backlog and a usage shortfall.
Within it, they resolve into a three-part diagnosis: a world that fell
silent at the chain's final link, since the act the course
assigned---analysis of the procedure---could be completed entirely
outside the world, and inside it nothing beyond watching was ever
asked; a flight that was never a malfunction, since the chain merely
shows each student following their greatest desire---both routes
served the governing end, and the greater desire lay with the more
convenient route; and a hail whose one unmatchable promise, the
meeting, went unconvened for want of a small reserve of shared time. Nothing in that diagnosis required
new concepts. The teaching clinic and Part II's frontier world differ
in every material and in no variable: an object of desire brought or
offered, an act housed or unhoused, a WE convened or left waiting. A
wholly authored world attunes without determining---and the collected
data has now shown both halves of that sentence at cohort scale, which
is what the dissertation carries forward from this section: not a
premise defended, but a finding held.
